\documentclass[10pt]{article}
\usepackage{amsmath}
\usepackage[english]{babel}
\usepackage[font=small,labelfont=bf]{caption}
\usepackage{geometry}
\usepackage[sort&compress]{natbib}
\usepackage{pxfonts}
\usepackage{graphicx}
\usepackage{setspace}
\usepackage{hyperref}
\usepackage{lineno}
\usepackage{xcolor}
%\usepackage[nolists, nomarkers, fighead]{endfloat}
%\renewcommand{\includegraphics}[2][]{}

\setcitestyle{notesep={; }}


\newcommand{\comment}[1]{}

\NewCommandCopy{\oldnameref}{\nameref}
\renewcommand{\nameref}[1]{\textit{\oldnameref{#1}}}

\doublespacing
\linenumbers

\title{TITLE PLACEHOLDER}

\author{Paxton C. Fitzpatrick\textsuperscript{1}, Alishba Tahir\textsuperscript{1, 2}, and Jeremy R. Manning\textsuperscript{1, *}\\
\vspace{-0.5em}\small{\textsuperscript{1}Department of Psychological and Brain
Sciences}\\
\small{Dartmouth College, Hanover, NH 03755, USA}\\
\vspace{-0.5em}\small{\textsuperscript{2}ALISHBA AFFILIATION PLACEHOLDER}\\
\small{ALISHBA AFFILIATION PLACEHOLDER}\\
\vspace{-0.5em}\small{\textsuperscript{*}Corresponding author:
Jeremy.R.Manning@Dartmouth.edu}}

\date{}

\begin{document}
\maketitle

\begin{abstract}\noindent ABSTRACT PLACEHOLDER
\end{abstract}


\section*{Introduction}


\section*{Results}


\section*{Discussion}


\section*{Methods}

Our experimental protocol was approved by the Committee for the Protection of Human Subjects at Dartmouth College.
We obtained written informed consent from all participants prior to the start of both experiment sessions.

\subsection*{Participants}

We enrolled a total of 70 Dartmouth undergraduate students in our study. 
Participants received optional course credit for enrolling and were pre-screened to ensure they had not previously seen any episode of either TV show viewed during the experiment. 
No statistical method was used to predetermine sample size. 
Prior to data processing and analysis, we excluded a total of 17 participants from our dataset: 4 encountered technical issues with the experiment, 2 self-reported difficulty with the task due to illness, 4 spoke for less than 10 minutes during the immediate recall phase, and 7 did not return for the second experiment session. 
The data we analyzed comprised the remaining 53 participants.

At the start of the experiment, we asked each participant to complete a demographic survey that included questions about their age, sex, race, and ethnicity.
Participants' ages ranged from 18 to 22 years (mean: 18.91 years; standard deviation: 1.01 years).
A total of 41 participants reported their sex as female and 12 participants reported their sex as male.
Participants reported their races as White (34 participants), Asian (16 participants), Black or African American (3 participants), American Indian or Alaska Native (1 participant), and other (1 participant).
(Note that some participants selected multiple racial categories.)
A total of 52 participants reported their ethnicity as ``Not Hispanic or Latino'' and 1 participant reported their ethnicity as ``Hispanic or Latino.''

We also asked participants their native spoken language; whether they had any hearing, speech, or color vision impairments; and about any current medications or recent injuries that could impact attention or memory.
All participants reported that their native spoken language was English, that they had no hearing or speech impairments, and that they had normal color vision.
A total of 52 participants reported no attention- or memory-relevant medications or injuries. One participant reported sustaining a concussion more than 3 years prior to the experiment and was not excluded from our dataset.

Additionally, before each of the two experiment sessions, we asked participants about their sleep, coffee consumption, and level of alertness. 
Participants reported having slept between 5 and 11 hours the night prior to session 1 (mean: 7.37 hours; standard deviation: 1.45 hours) and between 5 and 11 hours the night prior to session 2 (mean: 7.39 hours; standard deviation: 1.23 hours). 
They reported having consumed between 0 and 3 cups of coffee on the day of session 1 (mean: 0.43 cups; standard deviation: 0.72 cups) and between 0 and 3 cups of coffee on the day of session 2 (mean: 0.48 cups; standard deviation: 0.68 cups).
Participants reported their current levels of alertness on a 5-point Likert scale. At the start of session 1, 1 participant reported feeling ``Very sluggish,'' 14 reported feeling ``A little sluggish,'' 18 reported feeling ``Neutral,'' 15 reported feeling ``A little alert,'' and 5 reported feeling ``Very alert.'' At the start of session 2, 3 participants reported feeling ``Very sluggish,'' 9 reported feeling ``A little sluggish,'' 23 reported feeling ``Neutral,'' 17 reported feeling ``A little alert,'' and 3 reported feeling ``Very alert.''


\textbf{@jeremy}: \textit{do we want to report class years and majors? And/or do we want to report post-experiment survey responses (how engaging did you find the episode, how easy/difficult was it to follow, how well do you think you recalled it, how well do you think you learned characters' names)?}


\subsection*{Experiment}\label{subsec:experiment}

The experiment took place over two sessions lasting approximately one hour each, separated by 6--8 days (mean: 6.76 days; standard deviation: 0.73 days).
Prior to the start of the experiment, participants were assigned to one of two groups (\textit{A} or \textit{B}), alternating based on enrollment order.
During session 1, all participants first viewed the pilot episode of the FX series \textit{Atlanta} (``The Big Bang''; duration: 24 min 26 s).
Immediately after viewing, participants were given up to one hour to recount what happened in the episode to the best of their abilities.
Following \cite{ChenEtal17}, participants were instructed to describe the episode in as much detail as possible and to try to speak for a minimum of 10 minutes.
They were told to try to recount the episode's events in their original order, but that they were free to go back and describe earlier events should they realize they had missed something.
After the immediate recall phase, participants were then given up to 5 minutes to predict what they thought would happen in the show's next episode.

Upon returning for session 2, all participants were again given up to one hour to recount the \textit{Atlanta} episode they had viewed the week prior.
Following the delayed recall phase, participants in Group \textit{A} then viewed the second episode of \textit{Atlanta} (``Streets on Lock''; duration: 21 min 57 s) while participants in Group \textit{B} instead viewed the pilot episode of the FOX series \textit{Arrested Development} (``Pilot''; duration: 20 min 37 s).
Both groups were then given up to one hour to recount the just-watched episode in the same manner and given the same instructions as in session 1. 


\subsection*{Analysis}

\subsubsection*{Statistics}

Wilcoxon tests were one-sided because our alternative hypothesis was directional (delayed matches own immediate \textit{better than} the episode, others' immediates, etc.).

\subsubsection*{Constructing episode and recall trajectories}\label{subsec:text-embedding}


\subsubsection*{Event-segmentation}\label{subsec:eventseg}


\subsubsection*{2D path embedding}\label{subsec:umap}


\subsubsection*{}\label{subsec:smoothness}


\subsubsection*{Creating knowledge and learning map visualizations}\label{subsec:knowledge-maps}


\section*{Data availability}

All of the data analyzed in this manuscript may be found at
https://github.com/Con\-text\-Lab/memory-dynamics.

\section*{Code availability}

All of the code for running our experiment and carrying out the analyses may be
found at https://github.com/Con\-text\-Lab/memory-dynamics.

\bibliographystyle{apa}
\bibliography{cdl}

\section*{Acknowledgements}

INDIVIDUAL ACKNOWLEDGEMENTS PLACEHOLDER. Our work was supported in part
by NSF CAREER Award Number 2145172 to J.R.M. The content is solely the
responsibility of the authors and does not necessarily represent the official
views of our supporting organizations. The funders had no role in study design,
data collection and analysis, decision to publish, or preparation of the
manuscript.


\section*{Author contributions}


\section*{Competing interests}

The authors declare no competing interests.


\end{document}
